\documentclass[10pt,a4paper]{amsart}
\usepackage[margin=19mm]{geometry}
\usepackage{amsmath,amssymb,mathtools}
\usepackage[scr=rsfs,cal=euler]{mathalfa}
\usepackage{xfrac}
\usepackage[unicode,hidelinks,bookmarks=false]{hyperref}
\newcommand{\ie}{i.e.\ }
\newcommand{\cf}{cf.\ }
\newcommand{\resp}{respectively\ }
\newcommand{\Subsection}{\subsection}
\providecommand{\nobreakdash}{\nobreak\hbox{-}}
\setlength{\emergencystretch}{2em}
\multlinegap0pt
\usepackage{amssymb}
\usepackage{bm}
\usepackage{mathtools}
\usepackage{xcolor}
\usepackage[shortlabels]{enumitem}
\newtheorem{proposition}{Proposition}
\usepackage{cleveref}
\crefname{proposition}{Proposition}{Propositions}
\newcommand{\bq}{\mathbb{Q}}
\newcommand{\dd}{\mathrm{d}}
\renewcommand{\geq}{\geqslant}
\title{A $q$-analogue of Mirzakhani's recursion for Weil--Petersson volumes}
\author{Norman Do, Paul Norbury}
\date{Source: arXiv:2510.12431v2 (CC BY 4.0)}
\begin{document}
\maketitle

\noindent\emph{Source and licence.} A $q$-analogue of Mirzakhani's recursion for Weil--Petersson volumes. Version arXiv:2510.12431v2 (CC BY 4.0), distributed by arXiv under the Creative Commons Attribution 4.0 International licence, http://creativecommons.org/licenses/by/4.0/, as recorded in the arXiv licence metadata for this exact version and shown on the abstract page https://arxiv.org/abs/2510.12431v2. Full licence notice: \texttt{licenses/arxiv-2510.12431-CC-BY-4.0.txt}.\par\smallskip

\noindent\textit{Editorial scope.} Counted unit: the article's proposition F2k+1 (source0.tex lines 477--486). The article's complete original proof of it is delivered with it (source0.tex lines 489--553, one \texttt{proof} environment of 65 lines). No mathematics is written, repaired or re-derived: the statement and the proof are the article's own lines, in the article's own notation, and no retained line is substituted or re-flowed.  Retained uncounted as the source-local context the proof needs: lines 313--315; lines 317--319.  Nothing was omitted from any retained span, so the package carries no omission record.  No retained line contains a citation, so the wrapper carries no bibliography and no bibliography entry.  No retained line contains a figure environment, an image-inclusion command or drawing code, so no image is delivered and the package declares no asset dependency.  Editorial formatting only, in addition: an AMS \texttt{amsart} preamble that restates the article's own declarations and macros used by the retained lines, namely the environments \texttt{proposition} on separate counters, together with the standard packages those lines need.  The retained environments print in this document as Proposition 1 in order of appearance, whereas the article prints the same items with the numbering of its own full document.  The complete native source of the article as distributed in the arXiv e-print for this exact version is bundled unchanged as \texttt{sources/187037/source0.tex}.
\begin{equation} \label{qident}
\sum_{m=1}^\infty(-1)^{m-1}q^{m(m-1)/2}(1+q^m)\frac{q^{mk}}{(1-q^m)^{2k}} = s_k \big( \zeta_q(2), \zeta_q(4), \ldots \big),
\end{equation}

\begin{equation} \label{eq:sk}
\exp \bigg(\sum_{m=1}^\infty \frac{p_m}{m} \, t^m \bigg) = \sum_{k=-\infty}^\infty s_k(p_1, p_2, \ldots) \, t^k.
\end{equation}

\begin{proposition} \label{F2k+1}
For each non-negative integer $k$, we have
\[
\frac{F_{2k+1}(y)}{(2k+1)!} = \int_0^\infty \!\! \frac{x^{2k+1}}{(2k+1)!} \, H_q(x,y) \, \dd x = \sum_{n=0}^{k+1} b_n \frac{y^{2k+2-2n}}{(2k+2-2n)!},
\]
where $b_0, b_1, b_2, \ldots \in \bq[\zeta_q(2), \zeta_q(4), \ldots]$ are defined by
\[
\sum_{n=0}^\infty b_nz^{2n} = \exp \bigg(\sum_{j=1}^\infty\frac{\zeta_q(2j)}{j} (4z^2)^j \bigg) = 1+4\zeta_q(2)z^2+8(\zeta_q(2)^2+\zeta_q(4))z^4 + \cdots.
\]
\end{proposition}

\begin{proof}
Integrability of $x^{2k+1}H_q(x,y)$ is a consequence of integrability of $x^{2k+1} \frac{\dd^{2n}}{\dd x^{2n}} h_q(x)$ for all $n \geq 0$, together with vanishing of the integrals for $n>k+1$.


We first prove integrability of $x^{2k+1} \frac{\dd^{2n}}{\dd x^{2n}} h_q(x)$ in the $n=0$ case. Consider the following inequality, which is satisfied for $0<q<1$ and any positive integer $m$.
\begin{equation} \label{qmcomp} 
q^{-m/2}-q^{m/2}=q^{-m/2}(1+q+q^2+ \cdots +q^{m-1})(1-q)> (q^{-m/2}+m-1)(1-q)>m(1-q)
\end{equation}
Here, the first inequality uses $q^{-m/2+k}+q^{m/2-k}>2$ for $k=1,2, \ldots, m-1$. A consequence of \cref{qmcomp} is the bound
\[
|h_q(x)| < \sum_{m=1}^\infty q^{m(m-1)/2}(1+q^m) \, e^{\frac{1}{2} x (q^{m/2}-q^{-m/2})} < 2\sum_{m=1}^\infty e^{-\frac{1}{2}xm(1-q)} = \frac{2}{e^{\frac{1}{2}(1-q)x}-1}.
\]
We know that $\frac{2x^{2k+1}}{e^{\frac{1}{2}(1-q)x}-1}$ is integrable for each $k\geq 0$. Furthermore, \cref{qmcomp} allows us to deduce pointwise convergence of the following partial sums, via upper bounds on the tail of the series.
\begin{multline*}
\lim_{N\to\infty} \sum_{m=1}^N (-1)^{m-1}q^{m(m-1)/2}(1+q^m) \, e^{\frac{x}{2(1-q)} (q^{m/2}-q^{-m/2})} \\
= \sum_{m=1}^\infty (-1)^{m-1}q^{m(m-1)/2}(1+q^m) \, e^{\frac{x}{2(1-q)} (q^{m/2}-q^{-m/2})}
\end{multline*}
By the dominated convergence theorem, $x^{2k+1}h_q(x)$ is Lebesgue integrable in $x$, and we can interchange the sum and integral to obtain the following for $0<q<1$.
\begin{align*} 
\int_0^\infty \!\! \frac{x^{2k+1}}{(2k+1)!} \, h_q(x) \, \dd x &= \sum_{m=1}^\infty (-1)^{m-1}q^{m(m-1)/2}(1+q^m) \int_0^\infty \!\! \frac{x^{2k+1}}{(2k+1)!} \, e^{\frac{1}{2}x (q^{m/2}-q^{-m/2})} \, \dd x \\
&=\sum_{m=1}^\infty (-1)^{m-1}q^{m(m-1)/2}(1+q^m)\frac{ 2^{2k+2}}{(q^{m/2}-q^{-m/2})^{2k+2}} \\
&=2^{2k+2} \, s_{k+1}\left(\zeta_q(2), \zeta_q(4), \ldots \right)
\end{align*} 
The last equality uses the $q$-series identity of \cref{qident}, where $s_k$ is the polynomial defined by \cref{eq:sk}, which leads to the definition of $b_n$ in the statement of the proposition.


When $n>0$, the proof of integrability of $x^{2k+1} \frac{\dd^{2n}}{\dd x^{2n}} h_q(x)$ is similar to the $n=0$ case.  Partition $h_q(x)=h_q^{(2n)}(x)+$ tail into the sum of the first $2n$ terms and the tail. Integrability of the first $2n$ terms of $x^{2k+1} \frac{\dd^{2n}}{\dd x^{2n}} h_q(x)$ again uses the inequality given in \cref{qmcomp}:
\begin{align*}
\frac{\dd^{2n}}{\dd x^{2n}} h_q^{(2n)}(x) 
&= \sum_{m=1}^{2n} q^{m(m-1)/2}(1+q^m)2^{-2n} \left(q^{m/2}-q^{-m/2}\right)^{2n}e^{\frac{x}{2} (q^{m/2}-q^{-m/2})} \\
&< 2^{1-2n}(1-q)^{2n}\sum_{m=1}^{2n} m^{2n}e^{-\frac{1}{2}xm(1-q)}
\end{align*}
The tail is dominated by an integrable function as follows.
\[
\sum_{m=2n+1}^\infty q^{m(m-1)/2}(1+q^m) (q^{m/2}-q^{-m/2})^{2n}e^{\frac{x}{2}\left(q^{m/2}-q^{-m/2}\right)} < 2\sum_{m=1}^\infty e^{-\frac{1}{2}xm(1-q)}=\frac{2}{e^{\frac{1}{2}(1-q)x}-1}
\]
The condition $m>2n$ ensures that $m(m-1)/2-mn\geq0$ so that $q^{m(m-1)/2}(1+q^m)\left(q^{m/2}-q^{-m/2}\right)^{2n}<2$. Hence, $\left| x^{2k+1} \frac{\dd^{2n}}{\dd x^{2n}} h_q(x) \right|$ is bounded above by an integrable function.
The same comparison on the tail proves pointwise convergence for the series. So as in the $n=0$ case, by the dominated convergence theorem, $x^{2k+1} \frac{\dd^{2n}}{\dd x^{2n}} h_q(x)$ is Lebesgue integrable in $x$ and we can interchange the sum and integral. The integral can be obtained from the $n=0$ case via integration by parts or it can be calculated as follows.
\begin{align*} 
& \int_0^\infty \!\! \frac{x^{2k+1}}{(2k+1)!} \frac{\dd^{2n}}{\dd x^{2n}} h_q(x) \, \dd x \\
={}& \sum_{m=1}^\infty (-1)^{m-1}q^{m(m-1)/2}(1+q^m)\tfrac{1}{2^{2n}} (q^{m/2}-q^{-m/2})^{2n} \int_0^\infty \!\! \frac{x^{2k+1}}{(2k+1)!}e^{\frac{x}{2}(q^{m/2}-q^{-m/2})} \, \dd x \\
={}& \sum_{m=1}^\infty (-1)^{m-1}q^{m(m-1)/2}(1+q^m)\frac{ 2^{2k+2-2n}}{(q^{m/2}-q^{-m/2})^{2k+2-2n}} \\
={}& 2^{2k+2-2n} \, s_{k+1-n}\left(\zeta_q(2), \zeta_q(4), \ldots \right).
\end{align*} 


When $n>k+1$, the $q$-series identity of \cref{qident} implies the vanishing
\[
\sum_{m=1}^\infty (-1)^{m-1}q^{m(m-1)/2}(1+q^m)(q^{m/2}-q^{-m/2})^{2n-2k-2} = 0.
\]


Putting this together, we obtain
\begin{align*}
\int_0^\infty \!\! \frac{x^{2k+1}}{(2k+1)!} \, H(x,y) \, \dd x %&= \frac{1}{2}\sum_{m=1}^\infty\sum_{n=0}^{k+1}(-1)^{m-1}q^{m(m-1)/2}(1+q^m) \frac{\left(q^{m/2}-q^{-m/2}\right)^{2n}}{(2n)! \, (q^{-m/2}-q^{m/2})^{2k+2}} \, y^{2n} \\
%&= \frac{1}{2} \sum_{i=0}^{k+1}\sum_{m=1}^\infty(-1)^{m-1}q^{m(m-1)/2}(1+q^m)\frac{q^{mi}}{(1-q^m)^{2i}}\frac{y^{2k+2-2i}}{(2k+2-2i)!} \\
%&
=\sum_{i=0}^{k+1} s_i ( \zeta_q(2), \zeta_q(4), \ldots) \, \frac{(2y)^{2k+2-2i}}{(2k+2-2i)!}.
\end{align*}
Equivalently, $F_{2k+1}(y)$ is a polynomial in $y$ given by
\[
\frac{F_{2k+1}(y)}{(2k+1)!}=\sum_{n=0}^{k+1} b_n\frac{y^{2k+2-2n}}{(2k+2-2n)!} 
\]
where the factor of $2^{2k+2-2n}$ is absorbed by the definition of $b_n$.
\end{proof}

\end{document}
